\documentclass[12pt,a4paper]{article}
\usepackage[utf8]{inputenc}
\usepackage{geometry}
\geometry{a4paper, margin=1in}
\usepackage{hyperref}
\usepackage{titlesec}
\usepackage{enumitem}

\title{\textbf{Comprehensive Project Progress Report}\\ \large CBL LU Sukkur Plant HSE Management System}
\author{Talha Waheed \and Roshni Baloch \and Nisha \and Ali Raza \and Shahzaib Hushain \and Agha Shabhi}
\date{August 2026}

\begin{document}

\maketitle

\section*{1. Executive Summary}
The Health, Safety, and Environment (HSE) Management System is a full-stack, enterprise-grade application developed to digitize, monitor, and enforce safety protocols at the CBL LU Sukkur Plant. The platform replaces manual safety logging with a real-time dashboard and centralized database, significantly reducing incident response times and enabling data-driven compliance analysis.

The project is currently in the \textbf{Late Development / Pre-Deployment Phase}, with core functionality (Dashboards, Auth, Module tracking) successfully built and active QA processes underway.

\section*{2. Team Structure \& Role Allocation}
The success of the HSE Management System relies on the specialized, collaborative efforts of the following team members:
\begin{itemize}[leftmargin=*]
    \item \textbf{Nisha (Project Lead \& Core Backend Developer):} Oversees the entire project trajectory. Owns the core backend architecture using Node.js/Express, API route definitions, service logic, and coordinates front-to-back integration.
    \item \textbf{Talha Waheed (Frontend Engineer \& Auth Specialist):} Leads the React/Vite Single Page Application (SPA) development. Engineered the dynamic, schema-driven UI and charts. Responsible for the Microsoft Entra SSO integration and JWT token authorization mechanisms bridging the frontend and backend.
    \item \textbf{Ali Raza (Backend Developer \& Database Architect):} Designed and modeled the relational database schema in MySQL using Sequelize ORM. Structured the foundational tables, foreign key relationships, and data seeding strategies.
    \item \textbf{Roshni Baloch (Frontend Designer \& UI Tester):} Translated functional requirements into the highly professional ``Stitch Design'' aesthetics. Ensured color harmony, responsive layouts, intuitive UX, and conducts initial frontend UI validation.
    \item \textbf{Shahzaib Hushain (Documentation Specialist):} Responsible for tracking the system's technical specifications, API documentation, and drafting user manuals for final handover.
    \item \textbf{Agha Shabhi (Quality Assurance \& Testing):} Handles end-to-end (E2E) testing, integration testing, and User Acceptance Testing (UAT) to guarantee system reliability prior to plant deployment.
\end{itemize}

\section*{3. Current Project Status \& Implemented Components}

\subsection*{3.1. Executive Dashboard (Frontend)}
\textbf{Developed By:} Talha Waheed (Logic) \& Roshni Baloch (Design) \\
\textbf{Technology Used:} React 19, TypeScript, Recharts, Tailwind CSS. \\
\textbf{Functionality:} A dynamic, data-rich control panel displaying real-time safety KPIs. Includes Leading/Lagging Indicator grids, CAPA (Corrective Action) donut charts, Incident Bar charts, and a 6-month safety trend line graph. Features a robust date and department filtering system. \\
\textbf{Purpose \& Importance:} This is the heart of the application for plant managers. It translates raw database entries into actionable visual insights, allowing leadership to instantly identify safety risks, track training compliance, and oversee incident trends without manually crunching numbers.

\subsection*{3.2. Dynamic Data Entry Modules (Frontend)}
\textbf{Developed By:} Talha Waheed \& Roshni Baloch \\
\textbf{Technology Used:} React 19, Vite, React Router. \\
\textbf{Functionality:} A unified, schema-driven form generation engine (\texttt{DataEntrySection.tsx}). Instead of hardcoding 8 different forms, the system uses JSON configurations to dynamically render forms for Hazard Reporting, Near Misses, Incident Logs, and Training Records. Includes intelligent auto-filling based on Employee IDs. \\
\textbf{Purpose \& Importance:} Drastically reduces frontend code duplication. It ensures that adding a new module in the future takes minutes instead of days. It provides a standardized, error-proof data entry experience for plant workers.

\subsection*{3.3. Relational Database Schema (Backend)}
\textbf{Developed By:} Ali Raza \\
\textbf{Technology Used:} MySQL 8, Sequelize ORM. \\
\textbf{Functionality:} A highly normalized database architecture. Features tables for Users, Employees, Departments, Hazards, Incidents, etc. Utilizes UUIDs for secure record identification and foreign keys to prevent data duplication. \\
\textbf{Purpose \& Importance:} Ensures absolute data integrity. If a department is renamed or an employee changes roles, the relational design means it only needs to be updated in one place. It prevents database anomalies and supports complex analytical queries.

\subsection*{3.4. Backend API Architecture \& Rate Limiting (Backend)}
\textbf{Developed By:} Nisha \\
\textbf{Technology Used:} Node.js, Express, Sequelize. \\
\textbf{Functionality:} Exposes RESTful endpoints for CRUD operations across all safety modules. Includes robust middleware, error handling, structured API responses, and rate limiting to prevent server overload. \\
\textbf{Purpose \& Importance:} Serves as the secure bridge between the database and the frontend. It enforces business rules (e.g., preventing invalid data from being saved) and ensures the application can scale to handle hundreds of concurrent users at the plant.

\subsection*{3.5. Microsoft SSO \& JWT Authentication (Full-Stack)}
\textbf{Developed By:} Talha Waheed \\
\textbf{Technology Used:} Microsoft Entra ID (Azure AD), JSON Web Tokens (JWT), React Context. \\
\textbf{Functionality:} Allows plant employees to log into the dashboard using their existing Microsoft corporate credentials. The backend generates secure JWT tokens which the frontend attaches to every subsequent API request to prove identity. \\
\textbf{Purpose \& Importance:} Critical for enterprise security. It eliminates the need for employees to remember separate passwords and allows the IT department to manage access centrally. It ensures that only authorized personnel can view or report sensitive safety data.

\section*{4. Next Remaining Phases \& Components}
While the core functionality is operational, the following phases remain to reach v1.0 deployment:

\subsection*{Phase 1: Authentication \& Authorization Finalization (Talha \& Nisha)}
\textbf{Goal:} While SSO is integrated, the system must enforce strict role checks (e.g., standard users can report hazards, but only "HSE Managers" can close CAPAs or delete records). \\
\textbf{Importance:} Prevents unauthorized data manipulation.

\subsection*{Phase 2: Comprehensive E2E Testing (Agha Shabhi \& Roshni Baloch)}
\textbf{Goal:} Agha Shabhi will simulate real-world plant scenarios, systematically breaking the app to find edge cases. Roshni will ensure cross-browser compatibility and UI responsiveness on plant floor tablets. \\
\textbf{Importance:} Safety applications demand zero downtime and high reliability. Testing guarantees the system won't crash during a critical incident report.

\subsection*{Phase 3: File \& Evidence Upload Integration (Nisha \& Talha)}
\textbf{Goal:} Allowing users to attach photos of hazards or incident evidence directly to their reports (via AWS S3 or local storage). \\
\textbf{Importance:} Visual evidence is a mandatory requirement for HSE compliance.

\subsection*{Phase 4: Final Documentation \& Handover (Shahzaib Hushain)}
\textbf{Goal:} Shahzaib will compile the technical architecture docs for future IT maintainers and create a simplified user guide for plant workers learning the new system. \\
\textbf{Importance:} Ensures the system's longevity and smooth adoption across the Sukkur plant.

\end{document}
